\scanpage{48}
20) \uline{Exemple} \deleted{\unsure{de dualité}} avec un corps de base, sep. clos p.ex.
\[
H^i_!(X,F)^\vee\simeq H^{-i}(X,DF)
\]
Si $F$ loc. ct.
\[
H^i_!(X,F)^\vee\simeq H^{2d-i}(X,\underline{\operatorname{Hom}}(F,\mu_{nX}^{\otimes d}))
\]
On retrouve le th. de finitude de Artin.

Si $S$ quelconque, on \unsure{trouve} si les $R^i_!f_*(F)$ sont loc \unsure{constant|constants} (p.ex. $f$ propre lisse!)
\[
\operatorname{Hom}(R^i_!f_*(F),(\mathbb Z/n\mathbb Z)_X)\simeq R^{-i}_{\deleted{\unsure{!}}}f_*(DF))
\]
(ce qui implique que les $R^{-i}_!f_*(DF)$ sont loc \unsure{constant|constants}) et plus spécialement, si $F$ est loc. constant
\[
\operatorname{Hom}(R^i_!f_*(F),(\mathbb Z/n\mathbb Z)_X)\simeq R^{2d-i}_{\deleted{\unsure{!}}}f_*(\operatorname{Hom}(F,\mu_{nX}^{\otimes d})
\]
Donc dans ces cas, on dispose de la \deleted{finitude}\inserted{\unsure{constructiblt.|constructibilité}} des $R^{\deleted{\unsure{i}}j}_{\deleted{\unsure{!}}}\deleted{\operatorname{Hom}}G$ $(G=DF\ldots$

On peut généraliser en supposant que l'hypothèse \unsure{n'est|est} satisfaite que pour $i\leq i_0$, la conclusion le \unsure{vaut|reste} pour $j\geq2d-i_0\ldots$
\mspage{47}
